\section{Experimental Details}\label{app:F}
Every experiment of Section~\ref{sec:experiments} is described here. Two further analyses, which Section~\ref{sec:experiments} does not use, are in Appendix~\ref{app:G} (identification floors of models whose degree exceeds the cap) and Appendix~\ref{app:pisa} (a real planned-missing archive).

\subsection{Data, Populations and Software}\label{app:capsens}
\paragraph{Data.}
WDBC (569 records, 30 continuous features) comes from scikit-learn \citep{wolberg1993}; Spambase (4,601 records) and Ionosphere (351 records) from the PMLB mirror of the UCI files \citep{spambase1999,ionosphere1989,romano2021pmlb}; the CDC Diabetes Health Indicators (253,680 respondents of the 2015 Behavioral Risk Factor Surveillance System) from a public mirror of the Kaggle release referenced by the UCI entry \citep{cdcdiabetes2017}; and the thyroid files (7,200 patients) from public mirrors of the UCI release \citep{quinlan1986thyroid}. Each mirror was checked by reproducing values computed from the original downloads to nine decimals, which also confirms the record order.

\paragraph{Populations.}
Every population $P$ is an empirical law of held-out or target records mixed with a uniform mass $\alpha$ over all profiles and both labels, $\alpha=1\%$ unless stated. Full support makes $W_\A(P)/n$ a local minimax risk (Proposition~\ref{prop:local}) and keeps every cell mass positive. Oracle designs are optimized at $P$; fitted designs are optimized on a reference law (training or source records, smoothed with $\alpha=10\%$) with a floor of $5\%/L$ on each of the $L$ interview probabilities of the final master problem, and evaluated at $P$. Every design is pointwise unbiased, so oracle and fitted designs differ only in variance.

\paragraph{Software.}
Python 3.13.5, NumPy 2.5.4, SciPy 1.18.1, scikit-learn 1.9.1, pandas 3.1.1, CVXPY 1.9.3 with Clarabel 0.11.1, single-threaded. All constants are computed by the solver of Appendix~\ref{app:engine}; it reproduces the values of an independent dense implementation on every overlapping task.

\subsection{Forest Updates}\label{app:forestexact}
\paragraph{Protocol.}
For each dataset and seed, a random permutation assigns $n_{\rm old}$ records to train the old forest, the next $n_{\rm new}$ to train the new forest, and the remaining $n_{\rm eval}$ records to the evaluation law: $(n_{\rm old},n_{\rm new},n_{\rm eval})=(100,200,269)$ for WDBC, $(500,1000,3101)$ for Spambase and $(60,120,171)$ for Ionosphere. Continuous features are binarized at their training medians; $m$ features are chosen at equally spaced positions of the eligible feature list, without labels. Each forest has twenty trees of depth $b$, $\sqrt m$ features per split and minimum leaf size $\max\{3,n_{\rm old}/30\}$. The contrast $d$ is the difference of the two forests' predicted probabilities on the cube; it has degree at most $b$. The cap is binding for $m>b$. At $b=3$ twenty splits per $(m,\text{dataset})$ cell are used; at $b=2$ and $b=4$ five splits. The cap $b=4$ was completed for $m=6$ only: the branching program over all depth-four trees with forty padded component trees did not finish for $m\ge9$ within our compute budget, and $b=5$ was not run. At $m=12$ the held-out records occupy a median of 171, 388 and 121 of the 4,096 profiles for WDBC, Spambase and Ionosphere (ranges 164--180, 373--405 and 113--130 over twenty splits). Spambase, whose 3,101 held-out records cover few profiles densely, gives a median exact ratio of 1.58 at $b=3$, $m=12$, close to WDBC (1.66) and Ionosphere (1.73), so the gains are not an artifact of sparse held-out sets; Appendix~\ref{app:capsens} also varies the smoothing mass and resamples the held-out records.

\paragraph{Designs and certificates.}
The fixed constant optimizes over all $\binom mb$ forms; the branching constant over all depth-$b$ trees by column generation. The restricted menu adds only the forty component trees to the forms; at $b\ge4$ their short paths are padded to depth $b$, a refinement within the cap, so the restricted-menu ratios at $b\ge4$ are slightly more favourable to branching than at $b\le3$. At $b=4$ column generation stops when the certified bounds agree to relative error $10^{-4}$ or no tree improves the master; over the 15 tasks the largest relative gap of the reported ratio is 2e-08, the median number of iterations 28 and the median time 40 seconds per task on one core. At $b\in\{2,3\}$ an independent dense implementation certified every task to $10^{-6}$.

\begin{table}[t]
\centering
\caption{Independent forest updates, cap $b=3$, 20 splits per cell. Exact fixed/branching variance ratio over all forms and all depth-3 trees (median and 10th--90th percentiles). Restr.: branching restricted to the 40 component trees and all triples. Fix/full: fixed-form variance relative to the full-information variance. Fitted: designs optimized on training records, evaluated held-out ($m=12$, five splits).}
\label{tab:forestexact}
\footnotesize
\setlength{\tabcolsep}{2.6pt}
\begin{tabular}{llrrrr}
\toprule
Data & $m$ & Exact ratio & Restr. & Fix/full & Fitted \\
\midrule
WDBC & 6 & 1.36 [1.27, 1.48] & 1.08 & 2.11 &  \\
 & 9 & 1.57 [1.43, 1.69] & 1.09 & 3.91 &  \\
 & 12 & 1.66 [1.53, 1.74] & 1.08 & 5.55 & 1.67 \\
Spambase & 6 & 1.41 [1.27, 1.52] & 1.03 & 2.31 &  \\
 & 9 & 1.45 [1.36, 1.59] & 1.03 & 3.09 &  \\
 & 12 & 1.58 [1.49, 1.66] & 1.04 & 4.75 & 1.60 \\
Ionosphere & 6 & 1.36 [1.20, 1.46] & 1.07 & 2.06 &  \\
 & 9 & 1.61 [1.45, 1.77] & 1.10 & 4.19 &  \\
 & 12 & 1.73 [1.58, 1.78] & 1.10 & 5.94 & 1.79 \\
\bottomrule
\end{tabular}
\end{table}

\begin{table}[h]
\centering
\caption{Forest updates by cap and dimension (medians over splits). Fixed/branch at $b=3$ uses twenty splits per dataset (Table~\ref{tab:forestexact}); the two-round columns use five splits per dataset and the solver of Appendix~\ref{app:engine}: $W_2/W_{\rm branch}$, the share of the log-gap that the second round recovers (median and minimum over tasks), and the largest relative gap between the certified bounds of the branching constant.}
\label{tab:capdetail}
\small
\begin{tabular}{rrrrrrrrrr}
\toprule
$b$ & $m$ & WDBC & Spambase & Ionosphere & fitted & $W_2/W_{\rm branch}$ & recovered & min & width \\
\midrule
2 & 6 & 1.31 & 1.35 & 1.31 & 1.31 & -- & -- & -- & -- \\
2 & 9 & 1.23 & 1.25 & 1.44 & 1.26 & -- & -- & -- & -- \\
2 & 12 & 1.27 & 1.35 & 1.42 & 1.35 & -- & -- & -- & -- \\
3 & 6 & 1.36 & 1.41 & 1.36 & -- & 1.023 & 92\% & 85\% & 6e-07 \\
3 & 9 & 1.57 & 1.45 & 1.61 & -- & 1.041 & 90\% & 86\% & 1e-04 \\
3 & 12 & 1.66 & 1.58 & 1.73 & 1.67 & 1.057 & 89\% & 86\% & 8e-05 \\
4 & 6 & 1.23 & 1.28 & 1.20 & 1.27 & 1.031 & 85\% & 80\% & 9e-05 \\
\bottomrule
\end{tabular}
\end{table}


\paragraph{Two-round designs.}
For five splits per dataset at $(m,b)\in\{(6,3),(9,3),(12,3),(6,4)\}$ the solver computes the optimal fixed, two-round (and at $b=4$ three-round) and branching designs with certified bounds (relative gap at most \nForestGapMax). Of the \nForestOverlapN{} tasks that the dense implementation had also solved, the fixed-to-branching ratios agree to a relative error of $2\times10^{-4}$ on \nForestOverlapAgree; on the remaining \nForestOverlapRest{} the retrained forests differ slightly from the stored ones ($\Delta$ differs in the third decimal) and the ratios differ by at most \nForestOverlapMax\%. The share of the log-gap that a class recovers is $1-\log(W_r/W_{\rm branch})/\log(W_{\rm fixed}/W_{\rm branch})$. Two rounds recover a median \nTwoForestMed\% (range \nTwoForestLo--\nTwoForestHi\%); at $b=4$ three rounds recover a median \nThreeForestMed\% (Table~\ref{tab:capdetail}).

\paragraph{Smoothing mass and resampling.}
For the fifteen $b=3$ tasks with $m=12$ (five splits per dataset) the exact constants are recomputed with $\alpha\in\{0.1\%,1\%,5\%,10\%\}$ and, at $\alpha=1\%$, on bootstrap resamples of the held-out records with the forests fixed.
\begin{center}\small\begin{tabular}{lrrrrr}\toprule
Variation & tasks & median ratio & median rel. change & max rel. change & median $|\log|$ change \\
\midrule
$\alpha=0.1\%$ & 15 & 1.654 & +0.0\% & 2.0\% & 0.003 \\
$\alpha=5\%$ & 15 & 1.660 & +0.8\% & 2.7\% & 0.008 \\
$\alpha=10\%$ & 15 & 1.690 & +1.8\% & 5.3\% & 0.018 \\
bootstrap of held-out records ($\alpha=1\%$) & 27 & 1.602 & -1.1\% & 7.5\% & 0.012 \\
\bottomrule\end{tabular}\end{center}


\paragraph{Screening for a gain from reference data.}
For every forest task at $b\in\{2,3,4\}$ and every screening run, the uncentered lower end $(\E_P[t^2]/\Gamma_P-\Delta^2)/\sigma_P^2$ of Appendix~\ref{app:population} is recomputed on the reference law alone and compared with the realized exact ratio.
\begin{center}\small\begin{tabular}{lrrrrr}\toprule
Tasks & $n$ & Spearman($\widehat L$, realized) & median realized$/\widehat L$ & realized$<\widehat L$ & rule: $\widehat L\le1.05$ $\Rightarrow$ realized $\le1.25$ \\
\midrule
forests & 240 & 0.52 & 1.08 & 75 & 2/2 \\
screening trees & 48 & 0.78 & 1.00 & 0 & 31/31 \\
\bottomrule\end{tabular}\end{center}

For a screening tree this quantity is a lower bound on the records ratio at the reference law. For a forest update it is not: no single tree computes the contrast, branching itself needs more records than complete data, and the quantity bounds $W_{\rm fixed}$ relative to the full-information variance; it exceeds the realized fixed-to-branching ratio in 75 of the 240 forest tasks. We therefore use it only as a screen. In these data the screen \emph{at most 1.05 $\Rightarrow$ do not plan for a gain} has no misses: all 17 screening runs with a realized gain above 1.25 have a reference value above 1.05, all 31 runs with a value at most 1.05 realize at most 1.25, and none of the 214 forest tasks with a realized gain above 1.25 has a value at most 1.05 (only two forest tasks have a value at most 1.05).

\subsection{Screening Questionnaires}\label{app:screening}\label{app:screendetail}
\paragraph{Items and targets.}
BRFSS items: high blood pressure, high cholesterol, body-mass index at least 30, smoking, stroke, heart disease, physical activity, heavy drinking, fair or poor general health, difficulty walking, age 60 or older, and income in the five lowest categories; the target is the diabetes-or-prediabetes indicator. Thyroid items: age above its upper quartile, sex, thyroxine treatment, thyroid surgery, query hypothyroid, goitre, TSH above its upper quartile, T3, TT4 and FTI below their lower quartiles, T4U above its upper quartile, and sickness, with quartiles from the training file and no labels; the target is any abnormal class.

\paragraph{Splits and trees.}
For the random-half target, a random half of the records is the source; for the shifted target, the source is one sex and the target the other. Up to 20,000 source records train a class-balanced CART tree of depth $b$ (minimum leaf size 50 for BRFSS, 10 for thyroid); up to 20,000 further source records are the reference for fitted designs. The estimand is the target accuracy of the tree's decision, $d\in\{-1,1\}$, so $M_{\rm branch}=\E[t^2]=1$ and branching attains the full-information variance. Seeds $0,1,2$ vary the split and the tree.

\paragraph{Designs.}
Best: the optimal fixed design over all $\binom{12}b$ forms. Fitted: the same program on the source reference. Label: the optimal fixed design whose form probabilities may depend on the label, $W^{\rm lab}_{\rm fixed}(P)$ of Proposition~\ref{prop:labeldesign}, computed as two certified programs, one per label. Strat.: the same with strata defined by the label and sex (random-half BRFSS targets, where both sexes occur). PS: the optimized path-set design of Proposition~\ref{prop:upper}, at the target law and fitted. H2: follow the root question, then use the best fixed form of the remaining size in the reached subcube. $r$-round (exact): the certified optimum over all $r$-round interviews. Two fit.: column generation over two-round interviews on the source reference with the $5\%$ floor, started from the oracle two-round design, evaluated at the target. U-p: the uniform design over the tree\textquotesingle s padded path sets with the tables that minimize the variance for that design. U-i: the uniform design over all $b$-subsets of the items the decision depends on (a single form if it depends on at most $b$ items), with the best tables for that design. U-HT: the uniform design over all forms with the monomial Horvitz--Thompson tables $g_T=\sum_{S\subseteq T}\widehat t(S)\chi_S/N(S)$, where $N(S)=\binom{m-|S|}{b-|S|}$ is the number of forms that contain $S$. U-opt: the uniform design with the tables that minimize the variance for that design, a weighted least-norm problem. $1/\rho_P^2$ and $N_P$: the ends of Theorem~\ref{thm:bracket}.
\begin{table}[h]
\centering
\caption{All screening settings (medians over three seeds). Ratios are records needed relative to branching. Best/Fitted: optimal and fitted fixed designs; Label, Strat.: optimal fixed designs whose form probabilities depend on the label, and on the label and sex; PS: optimized path-set design (oracle/fitted); H2: heuristic two-round design; $r{=}2,3$: exact optimal $r$-round designs; 2 fit.: fitted two-round design; U-p, U-i: uniform designs over the tree\textquotesingle s path sets and over the $b$-subsets of the items the rule depends on, with optimal scores; U-HT, U-opt: uniform design over all forms with monomial Horvitz--Thompson scores and with optimal scores; $1/\rho_P^2$, $N_P$: ends of Theorem~\ref{thm:bracket}.}
\label{tab:screeningdetail}
\scriptsize
\setlength{\tabcolsep}{1.7pt}
\begin{tabular}{lllrrrrrrrrrrrrrrrrrr}
\toprule
Data & $b$ & Target & Acc. & Paths & Best & Fitted & Label & Strat. & PS & PS fit & H2 & $r{=}2$ & $r{=}3$ & 2 fit. & U-p & U-i & U-HT & U-opt & $1/\rho_P^2$ & $N_P$ \\
\midrule
BRFSS & 3 & random half & 0.629 & 2 & 1.00 & 1.04 & 1.00 & 1.00 & 1.00 & 1.00 & 1.00 & 1.00 & -- & 1.00 & 1.06 & 1.00 & 20.4 & 11.2 & 1.00 & 1.5 \\
BRFSS & 3 & other sex & 0.741 & 1 & 1.00 & 1.04 & 1.00 & -- & 1.00 & 1.00 & 1.00 & 1.00 & -- & 1.00 & 1.00 & 1.00 & 44.6 & 23.8 & 1.00 & 1.0 \\
BRFSS & 4 & random half & 0.664 & 4 & 1.46 & 1.52 & 1.45 & 1.45 & 1.50 & 1.52 & 1.14 & 1.00 & 1.00 & 1.01 & 1.77 & 2.30 & 43.2 & 21.5 & 1.18 & 3.2 \\
BRFSS & 4 & other sex & 0.681 & 4 & 1.52 & 1.61 & 1.51 & -- & 1.57 & 1.62 & 1.13 & 1.00 & 1.00 & 1.01 & 1.97 & 2.29 & 47.6 & 22.6 & 1.19 & 3.7 \\
BRFSS & 5 & random half & 0.675 & 7 & 1.28 & 1.32 & 1.28 & 1.27 & 1.28 & 1.29 & 1.15 & 1.02 & 1.00 & 1.02 & 1.66 & 1.89 & 26.9 & 11.9 & 1.04 & 5.2 \\
BRFSS & 5 & other sex & 0.734 & 6 & 1.43 & 1.52 & 1.43 & -- & 1.44 & 1.50 & 1.27 & 1.00 & 1.00 & 1.01 & 1.88 & 2.06 & 30.5 & 14.5 & 1.13 & 5.5 \\
BRFSS & 6 & random half & 0.719 & 14 & 1.53 & 1.58 & -- & -- & 1.53 & 1.56 & 1.40 & 1.06 & -- & -- & 2.41 & 3.09 & 29.7 & 11.6 & 1.13 & 11.0 \\
BRFSS & 6 & other sex & 0.722 & 13 & 1.60 & 1.70 & -- & -- & 1.65 & 1.70 & 1.39 & 1.01 & 1.00 & -- & 2.21 & 4.30 & 25.3 & 10.2 & 1.13 & 9.8 \\
Thyroid & 3 & random half & 0.823 & 3 & 1.00 & 1.04 & 1.00 & -- & 1.00 & 1.00 & 1.00 & 1.00 & -- & 1.04 & 1.19 & 1.00 & 32.4 & 10.8 & 1.00 & 2.4 \\
Thyroid & 3 & other sex & 0.819 & 3 & 1.00 & 1.04 & 1.00 & -- & 1.00 & 1.01 & 1.00 & 1.00 & -- & 1.00 & 1.37 & 1.00 & 31.9 & 10.7 & 1.00 & 2.8 \\
Thyroid & 4 & random half & 0.823 & 3 & 1.00 & 1.04 & 1.00 & -- & 1.00 & 1.00 & 1.00 & 1.00 & 1.00 & 1.04 & 1.15 & 1.00 & 11.3 & 4.9 & 1.00 & 2.4 \\
Thyroid & 4 & other sex & 0.848 & 4 & 1.00 & 1.04 & 1.00 & -- & 1.00 & 1.01 & 1.00 & 1.00 & 1.00 & 1.00 & 1.60 & 1.00 & 30.3 & 10.4 & 1.00 & 3.3 \\
Thyroid & 5 & random half & 0.823 & 6 & 1.00 & 1.03 & -- & -- & 1.00 & 1.00 & 1.00 & 1.00 & 1.00 & -- & 1.11 & 1.00 & 6.1 & 3.2 & 1.00 & 5.5 \\
Thyroid & 5 & other sex & 0.848 & 5 & 1.00 & 1.04 & -- & -- & 1.00 & 1.01 & 1.00 & 1.00 & 1.00 & -- & 1.27 & 1.00 & 10.7 & 4.9 & 1.00 & 4.0 \\
Thyroid & 6 & random half & 0.823 & 7 & 1.00 & 1.03 & -- & -- & 1.00 & 1.00 & 1.00 & 1.00 & 1.00 & -- & 1.09 & 1.00 & 4.0 & 2.4 & 1.00 & 6.1 \\
Thyroid & 6 & other sex & 0.848 & 5 & 1.00 & 1.04 & -- & -- & 1.00 & 1.01 & 1.00 & 1.00 & 1.00 & -- & 1.27 & 1.00 & 5.7 & 3.1 & 1.00 & 4.0 \\
\bottomrule
\end{tabular}
\end{table}

\begin{table}[h]
\centering
\caption{Every BRFSS screening run (twelve items; R: random-half target, S: other-sex target; seeds 0--2) and every 21-item run: records needed relative to branching. Items: number of items the decision depends on. Two: exact optimal two-round design; a bracket gives the certified lower and upper bounds reached within the time budget. -- not computed.}
\label{tab:runs}
\scriptsize
\setlength{\tabcolsep}{3.2pt}
\begin{tabular}{llrrrrrrrrrr}
\toprule
Run & $b$ & Acc. & Items & Paths & Best & Fitted & Label & Two & Two fit. & PS & $1/\rho_P^2$ \\
\midrule
R0 & 3 & 0.604 & 2 & 2 & 1.000 & 1.035 & 1.000 & 1.000 & 1.035 & 1.000 & 1.000 \\
R1 & 3 & 0.736 & 3 & 4 & 1.000 & 1.040 & 1.000 & 1.000 & 1.000 & 1.000 & 1.000 \\
R2 & 3 & 0.629 & 3 & 2 & 1.000 & 1.037 & 1.000 & 1.000 & 1.000 & 1.000 & 1.000 \\
S0 & 3 & 0.742 & 3 & 1 & 1.000 & 1.040 & 1.000 & 1.000 & 1.000 & 1.000 & 1.000 \\
S1 & 3 & 0.720 & 3 & 1 & 1.000 & 1.040 & 1.000 & 1.000 & 1.000 & 1.000 & 1.000 \\
S2 & 3 & 0.741 & 4 & 2 & 1.338 & 1.409 & 1.338 & 1.000 & 1.008 & 1.338 & 1.070 \\
R0 & 4 & 0.664 & 6 & 4 & 1.457 & 1.524 & 1.451 & 1.000 & 1.009 & 1.496 & 1.153 \\
R1 & 4 & 0.659 & 5 & 4 & 1.406 & 1.461 & 1.404 & 1.000 & 1.009 & 1.458 & 1.176 \\
R2 & 4 & 0.687 & 6 & 3 & 1.550 & 1.611 & 1.543 & 1.000 & 1.010 & 1.608 & 1.192 \\
S0 & 4 & 0.664 & 4 & 2 & 1.000 & 1.035 & 1.000 & 1.000 & 1.035 & 1.000 & 1.000 \\
S1 & 4 & 0.684 & 7 & 5 & 1.749 & 1.867 & 1.740 & 1.000 & 1.011 & 1.817 & 1.213 \\
S2 & 4 & 0.681 & 6 & 4 & 1.517 & 1.610 & 1.511 & 1.000 & 1.006 & 1.570 & 1.188 \\
R0 & 5 & 0.675 & 6 & 7 & 1.184 & 1.231 & 1.183 & 1.018 & 1.020 & 1.189 & 1.031 \\
R1 & 5 & 0.678 & 7 & 8 & 1.279 & 1.323 & 1.275 & 1.000 & 1.007 & 1.281 & 1.042 \\
R2 & 5 & 0.663 & 8 & 5 & 1.513 & 1.581 & 1.508 & 1.040 & 1.046 & 1.544 & 1.099 \\
S0 & 5 & 0.692 & 7 & 6 & 1.260 & 1.322 & 1.259 & 1.000 & 1.006 & 1.261 & 1.070 \\
S1 & 5 & 0.746 & 8 & 6 & 1.537 & 1.645 & 1.535 & 1.000 & 1.011 & 1.564 & 1.144 \\
S2 & 5 & 0.734 & 7 & 9 & 1.430 & 1.522 & 1.425 & 1.000 & 1.011 & 1.439 & 1.132 \\
R0 & 6 & 0.729 & 10 & 14 & 1.982 & 2.059 & -- & [1.000, 1.320] & -- & 2.022 & 1.264 \\
R1 & 6 & 0.719 & 9 & 17 & 1.530 & 1.580 & -- & [1.000, 1.063] & -- & 1.534 & 1.132 \\
R2 & 6 & 0.686 & 9 & 13 & 1.417 & 1.475 & -- & [1.012, 1.031] & -- & 1.438 & 1.066 \\
S0 & 6 & 0.690 & 10 & 12 & 1.604 & 1.700 & -- & 1.000 & -- & 1.645 & 1.119 \\
S1 & 6 & 0.725 & 10 & 13 & 1.714 & 1.827 & -- & 1.013 & -- & 1.767 & 1.150 \\
S2 & 6 & 0.722 & 10 & 20 & 1.561 & 1.664 & -- & [1.000, 1.109] & -- & 1.570 & 1.132 \\
\addlinespace[2pt]
21 items, R0 & 3 & 0.604 & 2 & 2 & 1.000 & 1.000 & -- & 1.000 & -- & 1.000 & 1.000 \\
21 items, R1 & 3 & 0.736 & 3 & 4 & 1.000 & 1.016 & -- & 1.000 & -- & 1.000 & 1.000 \\
21 items, R2 & 3 & 0.629 & 3 & 2 & 1.000 & 1.002 & -- & 1.000 & -- & 1.000 & 1.000 \\
21 items, R0 & 4 & 0.664 & 6 & 5 & 1.457 & 1.535 & -- & 1.000 & -- & 1.496 & 1.153 \\
21 items, R1 & 4 & 0.659 & 5 & 4 & 1.406 & 1.471 & -- & 1.000 & -- & 1.458 & 1.176 \\
21 items, R2 & 4 & 0.687 & 6 & 3 & 1.550 & 1.623 & -- & 1.000 & -- & 1.608 & 1.192 \\
\bottomrule
\end{tabular}
\end{table}


\paragraph{Label-dependent and stratified forms.}
Over the \nLabelN{} runs, letting the form probabilities depend on the label lowers the records ratio by at most \nLabelMax\% (certified gap between the bounds at most \nLabelCert), and stratifying in addition by sex by at most \nStratMax\% over \nStratN{} runs. The prevalence of the target is about 14\%, so the two labels do weight the tree's paths differently, but the forms that serve one label serve the other almost as well.

\paragraph{Two-round designs.}
A two-round interview that asks a first batch $B_1$ computes the decision exactly if and only if, for every answer pattern on $B_1$, the restricted decision depends on at most $b-|B_1|$ further items. If such a batch exists, the interview observes $t$ and $W_2=W_{\rm branch}$; at a full-support law the converse holds as well, because a design with variance $\sigma_P^2$ must return $t$ at every profile. We enumerate all first batches inside the items the decision depends on. Of the \nBrfssRuns{} BRFSS runs, \nCompBrfss{} are two-round computable, and the enumeration agrees with the solver on all \nCompAgree{} runs for which the solver also computed the two-round optimum (ratio one exactly when computable). The trees that are not two-round computable are the larger ones: two of the six at $b=5$ and five of the six at $b=6$; class-balanced CART produces unbalanced trees with few distinct paths at the smaller caps. Exact two-round optima were computed by the solver for every run at $b\le5$. At $b=6$ one run is two-round computable. For the other five, column generation was given one hour per run on one core (about fifty to sixty iterations and five to seven hundred generated interviews). One run closed its bracket to a relative gap of $1.3\times10^{-4}$ and is counted as solved; for the other \nTwoBounded{} Table~\ref{tab:runs} reports the bracket reached. Its upper end is the variance of a feasible, exactly unbiased two-round design. Its lower end is one where the dual bound of the stopped run was not yet above the full-information bound; for these runs the ratio is nevertheless known to exceed one, since the trees are not two-round computable. The forest updates are not two-round computable either (all sixty exact two-round ratios exceed one), and neither is the pointer tree.

\subsection{All Twenty-One BRFSS Variables}\label{app:scale}
All 21 variables of the CDC file are binarized at pre-specified cut points: high blood pressure, high cholesterol, cholesterol check, body-mass index at least 30, smoking, stroke, heart disease, physical activity, fruit, vegetables, heavy drinking, any health coverage, cost barrier to a doctor visit, fair or poor general health, at least 14 days of poor mental health, at least 14 days of poor physical health, difficulty walking, sex, age 60 or older, no college education, and income in the five lowest categories. Splits, trees and references are as for the twelve-item study (random-half target). The form program has $\binom{21}b$ forms, 1,330 at $b=3$ and 5,985 at $b=4$, and the profile cube has $2^{21}$ elements; the solver never enumerates the cube, and at $b=4$ one master problem over all forms certifies the optimum. In all three seeds at $b=4$ the tree selects only items of the twelve-item study and coincides with the tree of that study (runs R0--R2 of Table~\ref{tab:runs}), so its constants coincide as well; the purpose of this study is the size of the program, whose optimum must be certified against all 5,985 forms. All three trees at $b=4$ are two-round computable, so their two-round ratio is exactly one. Label-dependent designs and the cap $b=5$ (20,349 forms) were not completed at $m=21$ within our compute budget.
\begin{table}[h]
\centering
\caption{All 21 BRFSS variables (random-half target; medians over three seeds). Forms: $\binom{21}b$. Time: median minutes for the certified fixed design on one core.}
\label{tab:scale}
\small
\begin{tabular}{rrrrrrrrrr}
\toprule
$b$ & forms & Acc. & items used & Best & Fitted & PS & $1/\rho_P^2$ & $N_P$ & time \\
\midrule
3 & 1330 & 0.629 & 4 & 1.00 & 1.00 & 1.00 & 1.00 & 1.5 & 0.0 \\
4 & 5985 & 0.664 & 6 & 1.46 & 1.54 & 1.50 & 1.18 & 3.5 & 6.1 \\
\bottomrule
\end{tabular}
\end{table}


\subsection{Learning the Design from a Pilot}\label{app:pilot}
Proposition~\ref{prop:selfcal} is asymptotic and discards the pilot. To measure the all-in cost of learning at finite sizes we use the BRFSS runs with a gain and the tree's path-set menu, whose optimum is within a few percent of the all-forms optimum (Table~\ref{tab:screeningdetail}). A pilot of $n_1$ target respondents is interviewed with the uniform design on the path sets. The pilot is split into two halves at random; the masses of the cells of each path-set form are estimated from one half, mixed with a uniform mass of $10\%$, and used to compute the tables for the uniform design that score the \emph{other} half. These scores are pointwise unbiased and independent of the records they score. The main sample of $N-n_1$ respondents uses the path-set design learned from the whole pilot with the $5\%$ floor. The final estimate averages all $N$ scores; its variance is $[n_1w_{\rm pilot}+(N-n_1)w_{\rm learned}]/N^2$ for the per-record variances $w$, which are computed exactly at the target law for 30 independent pilots per run and size.
\begin{table}[h]
\centering
\caption{Learning the design from a pilot of the audit itself, with the pilot charged (BRFSS runs with a gain; medians over runs of the median over 30 pilots). Pilot: per-record variance of the cross-fitted scores of the $n_1$ pilot respondents (uniform design on the path sets). Learned: per-record variance of the path-set design learned from the pilot. All-in: variance of the estimator that uses all $N=4n_1$ records with weights proportional to sample sizes, $[n_1\cdot\text{pilot}+3n_1\cdot\text{learned}]/N$. All entries are relative to branching; the last columns give the oracle path-set design and the oracle design over all forms. Cells: number of (form, pattern, label) masses to estimate.}
\label{tab:pilot}
\small
\begin{tabular}{rrrrrrrrr}
\toprule
$b$ & runs & cells & $n_1$ & pilot & learned & all-in & oracle PS & oracle all forms \\
\midrule
3 & 1 & 32 & 250 & 1.56 & 1.37 & 1.42 & 1.34 & 1.34 \\
 &  &  & 500 & 1.53 & 1.36 & 1.40 &  &  \\
 &  &  & 1000 & 1.52 & 1.35 & 1.39 &  &  \\
 &  &  & 2000 & 1.51 & 1.35 & 1.39 &  &  \\
4 & 5 & 128 & 250 & 2.43 & 1.80 & 1.96 & 1.57 & 1.52 \\
 &  &  & 500 & 2.13 & 1.71 & 1.75 &  &  \\
 &  &  & 1000 & 1.99 & 1.64 & 1.67 &  &  \\
 &  &  & 2000 & 1.90 & 1.61 & 1.65 &  &  \\
5 & 6 & 416 & 250 & 3.26 & 2.46 & 2.61 & 1.36 & 1.35 \\
 &  &  & 500 & 2.95 & 1.90 & 2.09 &  &  \\
 &  &  & 1000 & 2.31 & 1.66 & 1.80 &  &  \\
 &  &  & 2000 & 1.99 & 1.54 & 1.63 &  &  \\
6 & 6 & 1728 & 250 & 4.53 & 4.15 & 4.24 & 1.61 & 1.58 \\
 &  &  & 500 & 5.57 & 4.21 & 4.59 &  &  \\
 &  &  & 1000 & 5.40 & 3.71 & 4.13 &  &  \\
 &  &  & 2000 & 4.79 & 2.76 & 3.30 &  &  \\
\bottomrule
\end{tabular}
\end{table}

With a quarter of the records in the pilot, the all-in price exceeds the oracle all-forms price by a median \nAllInOverFourc\% and \nAllInOverFivec\% at $b=4,5$ for $N=4{,}000$ and by \nAllInOverFourd\% and \nAllInOverFived\% for $N=8{,}000$. The excess falls roughly in proportion to the number of cells per pilot respondent; at $b=6$ pilots of this size are too small. A design fitted on a large reference does better at every cap (Table~\ref{tab:runs}), so a pilot is the tool for the case in which no reference exists, and the all-forms menu, with $\binom{12}b2^{b+1}$ cells, is out of reach of pilots of this size. A non-asymptotic guarantee for the two-phase procedure is not proved here.

\subsection{One Set of Records for Several Rules}\label{app:reuse}
An optimized design concentrates on the forms that serve one contrast. To measure what this costs when the records are later used for another rule, we take the other-sex BRFSS target, whose population is the same for every seed, and treat the three trees of seeds 0--2 at each cap as three rules. For every ordered pair we compute the best pointwise-unbiased scores for the second rule on records collected under the first rule's optimal design, mixed with a $5\%$ uniform floor over all forms so that every contrast of degree at most $b$ stays identifiable. For every pair and for the triple we also solve the fixed-form program for several contrasts of Appendix~\ref{app:population}, which minimizes the sum of the rules' records ratios over one shared form distribution.
\begin{table}[h]
\centering
\caption{Using one set of records for several rules (BRFSS, other-sex target; the three trees of seeds 0--2 at each cap are rules A, B, C). Entries are records needed relative to branching for that rule alone. Own: the rule\textquotesingle s own optimal fixed design. Reuse: best scores for the rule on records collected under another rule\textquotesingle s optimal design mixed with a 5\% uniform floor (minimum and maximum over the two other designs). Joint: the fixed design that minimizes the sum of the records ratios of the rules in the group.}
\label{tab:reuse}
\small
\begin{tabular}{rlrrrrrrr}
\toprule
$b$ & Rule & Own & Reuse min & Reuse max & Joint A,B & Joint A,C & Joint B,C & Joint A,B,C \\
\midrule
4 & A & 1.00 & 1.97 & 2.14 & 1.40 & 1.34 &  & 1.51 \\
 & B & 1.75 & 30.6 & 316.7 & 1.94 &  & 1.76 & 1.87 \\
 & C & 1.52 & 1.69 & 286.6 &  & 1.67 & 1.60 & 1.67 \\
5 & A & 1.26 & 20.4 & 24.5 & 1.36 & 1.40 &  & 1.43 \\
 & B & 1.54 & 10.2 & 42.2 & 1.59 &  & 1.54 & 1.56 \\
 & C & 1.43 & 1.55 & 41.2 &  & 1.59 & 1.48 & 1.58 \\
\bottomrule
\end{tabular}
\end{table}

Reuse can be nearly free or nearly useless: the reused records need between \nReuseMin{} and \nReuseMax{} times the records of branching, depending on how many of the second rule's paths the first design happens to cover. A joint design for rules known in advance costs each rule a median \nJointPenMed\% (at most \nJointPenMax\%) more than its own optimum. For comparison, giving each of three rules its own third of the respondents and administering its tree costs each rule a factor of three.

\subsection{Model-Based Alternatives}\label{app:regimp}\label{app:shift}
\paragraph{Simulated audits.}
Each replication draws $n=1000$ target respondents with replacement. Branching averages $t$. The fixed-form score draws a form from the fitted design and averages its score. Truth is the contrast on the full target sample. The trees and fitted designs are those of seed 0, so the random-half and sex settings are runs R0 and S0 of Table~\ref{tab:runs}; in particular the exact records ratio of the sex-shift tree at $b=4$ is 1.000 (fitted 1.006), so the fitted form and branching have nearly equal standard deviations in the simulation, whereas Table~\ref{tab:screeningdetail} reports medians over three seeds.

\paragraph{Imputation from complete source records.}
Regression imputation uses the same observed forms and answers and replaces $t$ by its conditional mean given the asked items and the label among the complete source records, with a naive standard error. Multiple imputation draws, for each of ten imputations, Bayesian-bootstrap weights over the source records, imputes the unasked items of each respondent from a donor that agrees on the asked items and the label, and combines by Rubin's rules. Both use information that the capped audit does not collect: complete records whose conditional law of unasked items is assumed to transport to the target. Before any result was computed we fixed four shifts defined by variables outside the twelve items: sex (female source, male target), education (college graduates, target without education beyond high school), cost barrier (no cost barrier to a doctor visit, target with a cost barrier) and physical-health burden (no poor physical-health days in the past month, target with at least 14). Source and target do not overlap in the shift variable, so an imputation model cannot condition on it. All four are reported in Table~\ref{tab:shift} (400 replications).
\begin{table}[h]
\centering
\caption{Imputation from complete source records: simulated BRFSS audits with $n=1000$ target respondents and 400 replications, for the random-half target (shift: none) and the four pre-specified shifts. The trees are those of seed 0; the first two blocks are runs R0 and S0 of Table~\ref{tab:runs}. Bias, SD and RMSE are multiplied by 100; coverage of nominal 95\% intervals has Monte Carlo standard error about 0.011. TV: total variation distance between source and target item profiles. Regression imputation: conditional mean of $t$ given the asked items and the label among the source records, naive standard error. Multiple imputation: ten proper hot-deck imputations from the same source records, Rubin\textquotesingle s rules.}
\label{tab:shift}
\scriptsize
\setlength{\tabcolsep}{2.6pt}
\begin{tabular}{lrrlrrrrr}
\toprule
Shift & $b$ & TV & Estimator & Bias & SD & RMSE & Coverage & Width \\
\midrule
none & 4 & 0 & branching & -0.14 & 2.96 & 2.96 & 0.955 & 11.70 \\
 &  &  & fitted fixed form & -0.11 & 3.60 & 3.60 & 0.965 & 14.42 \\
 &  &  & regression imputation & -0.45 & 2.59 & 2.63 & 0.955 & 10.27 \\
 &  &  & multiple imputation & -0.42 & 2.58 & 2.61 & 0.993 & 13.21 \\
none & 5 & 0 & branching & -0.05 & 2.92 & 2.92 & 0.950 & 11.60 \\
 &  &  & fitted fixed form & -0.06 & 3.19 & 3.18 & 0.955 & 12.85 \\
 &  &  & regression imputation & -0.06 & 2.67 & 2.67 & 0.958 & 10.84 \\
 &  &  & multiple imputation & -0.06 & 2.68 & 2.68 & 0.975 & 12.41 \\
sex & 4 & 0.21 & branching & 0.05 & 3.06 & 3.05 & 0.953 & 11.69 \\
 &  &  & fitted fixed form & 0.06 & 3.12 & 3.11 & 0.948 & 11.89 \\
 &  &  & regression imputation & 0.32 & 3.01 & 3.03 & 0.940 & 11.50 \\
 &  &  & multiple imputation & 0.33 & 3.02 & 3.03 & 0.955 & 11.87 \\
sex & 5 & 0.21 & branching & 0.08 & 2.89 & 2.89 & 0.965 & 11.42 \\
 &  &  & fitted fixed form & 0.11 & 3.44 & 3.44 & 0.938 & 13.12 \\
 &  &  & regression imputation & -0.45 & 2.80 & 2.84 & 0.948 & 10.72 \\
 &  &  & multiple imputation & -0.47 & 2.81 & 2.85 & 0.983 & 12.22 \\
education & 4 & 0.47 & branching & 0.08 & 3.12 & 3.12 & 0.948 & 12.28 \\
 &  &  & fitted fixed form & -0.06 & 3.60 & 3.60 & 0.960 & 14.20 \\
 &  &  & regression imputation & 2.39 & 2.79 & 3.67 & 0.863 & 11.28 \\
 &  &  & multiple imputation & 2.38 & 2.83 & 3.69 & 0.920 & 13.27 \\
education & 5 & 0.47 & branching & 0.04 & 3.11 & 3.11 & 0.950 & 12.25 \\
 &  &  & fitted fixed form & -0.03 & 3.39 & 3.38 & 0.953 & 13.28 \\
 &  &  & regression imputation & 1.67 & 2.96 & 3.39 & 0.902 & 11.76 \\
 &  &  & multiple imputation & 1.66 & 2.97 & 3.39 & 0.935 & 12.73 \\
cost barrier & 4 & 0.41 & branching & -0.09 & 3.23 & 3.22 & 0.945 & 12.01 \\
 &  &  & fitted fixed form & -0.18 & 3.80 & 3.80 & 0.943 & 14.75 \\
 &  &  & regression imputation & 0.39 & 2.82 & 2.85 & 0.930 & 10.76 \\
 &  &  & multiple imputation & 0.39 & 2.84 & 2.86 & 0.980 & 13.31 \\
cost barrier & 5 & 0.41 & branching & 0.01 & 3.13 & 3.12 & 0.935 & 11.75 \\
 &  &  & fitted fixed form & -0.10 & 3.89 & 3.89 & 0.938 & 14.45 \\
 &  &  & regression imputation & -0.72 & 2.84 & 2.93 & 0.925 & 10.58 \\
 &  &  & multiple imputation & -0.70 & 2.92 & 3.00 & 0.960 & 13.01 \\
poor health & 4 & 0.68 & branching & 0.21 & 3.17 & 3.17 & 0.950 & 12.39 \\
 &  &  & fitted fixed form & 0.23 & 3.23 & 3.23 & 0.950 & 12.59 \\
 &  &  & regression imputation & 0.65 & 3.15 & 3.21 & 0.960 & 12.23 \\
 &  &  & multiple imputation & 0.65 & 3.15 & 3.22 & 0.965 & 12.57 \\
poor health & 5 & 0.68 & branching & 0.09 & 3.17 & 3.16 & 0.945 & 12.37 \\
 &  &  & fitted fixed form & 0.14 & 3.94 & 3.94 & 0.940 & 15.04 \\
 &  &  & regression imputation & 5.49 & 2.96 & 6.24 & 0.515 & 11.44 \\
 &  &  & multiple imputation & 5.48 & 2.96 & 6.23 & 0.667 & 13.23 \\
\bottomrule
\end{tabular}
\end{table}


Under the random-half target, MI intervals over-cover (0.975--0.993), so Rubin\textquotesingle s variance is conservative there. For a frozen imputation bias $\mathfrak b_I$ and per-record variances $\sigma_I^2$ and $\sigma_U^2$ of the imputation and design-based estimators, imputation has the smaller mean squared error exactly when $n\mathfrak b_I^2<\sigma_U^2-\sigma_I^2$, so a biased imputation can win at small $n$ and lose as the audit grows.

\paragraph{A working model on the capped records.}
Planned-missing data are usually analysed by maximum likelihood or multiple imputation under a working model for the items, fitted to the incomplete records themselves; missingness is by design, so no transport assumption is involved. We fit the log-linear model with all main effects and two-way interactions among the label and the items that the screening rule uses (items the rule ignores carry no information about the target) by EM, with one iterative-proportional-fitting sweep per M-step, and estimate $\Delta$ by $\sum_vp(v)t(v)$. Intervals use a delete-a-group jackknife with eight groups and a Student quantile. The estimator is applied to the records of the fitted fixed-form design and to records of balanced random forms (uniform over all $\binom{12}b$ forms), in the same simulated audits (\nEmReps{} replications).
\begin{table}[h]
\centering
\caption{A working model on the capped records: simulated BRFSS audits with $n=1000$ and $n=16,000$ target respondents (400 replications each). Design-based scores (branching; fitted fixed form) against maximum likelihood under a two-way log-linear model for the label and the items the rule uses, fitted by EM to the fitted-form records (EM, fitted forms) and to balanced random forms (EM, uniform forms). Bias, SD and RMSE are multiplied by 100; coverage of nominal 95\% intervals (delete-a-group jackknife for EM) has Monte Carlo standard error about 0.011. -- not computed.}
\label{tab:em}
\scriptsize
\setlength{\tabcolsep}{2.8pt}
\begin{tabular}{lrrlrrrrrrrr}
\toprule
 & & & & \multicolumn{4}{c}{$n=1000$} & \multicolumn{4}{c}{$n=16,000$} \\
Shift & $b$ & items & Estimator & Bias & SD & RMSE & Cov. & Bias & SD & RMSE & Cov. \\
\midrule
none & 4 & 6 & branching & -0.01 & 3.05 & 3.05 & 0.950 & -0.01 & 0.76 & 0.76 & 0.943 \\
 &  &  & fitted fixed form & -0.07 & 3.71 & 3.70 & 0.950 & -0.03 & 0.88 & 0.88 & 0.953 \\
 &  &  & EM, fitted forms & 0.74 & 2.88 & 2.97 & 0.935 & 0.77 & 0.76 & 1.08 & 0.855 \\
 &  &  & EM, uniform forms & 0.20 & 4.91 & 4.91 & 0.953 & 0.50 & 1.19 & 1.29 & 0.922 \\
none & 5 & 6 & branching & 0.00 & 3.06 & 3.06 & 0.940 & -0.03 & 0.75 & 0.75 & 0.940 \\
 &  &  & fitted fixed form & 0.04 & 3.37 & 3.37 & 0.950 & -0.03 & 0.84 & 0.84 & 0.938 \\
 &  &  & EM, fitted forms & 0.20 & 2.75 & 2.75 & 0.948 & 0.14 & 0.71 & 0.72 & 0.932 \\
 &  &  & EM, uniform forms & 0.15 & 4.25 & 4.25 & 0.960 & 0.18 & 1.13 & 1.14 & 0.932 \\
sex & 4 & 4 & branching & 0.15 & 2.96 & 2.96 & 0.958 & -0.05 & 0.77 & 0.77 & 0.935 \\
 &  &  & fitted fixed form & 0.18 & 3.00 & 3.00 & 0.955 & -0.04 & 0.78 & 0.78 & 0.943 \\
 &  &  & EM, fitted forms & 0.51 & 2.79 & 2.84 & 0.955 & 0.37 & 0.69 & 0.78 & 0.920 \\
 &  &  & EM, uniform forms & -0.11 & 5.33 & 5.33 & 0.925 & 0.35 & 1.22 & 1.27 & 0.943 \\
sex & 5 & 7 & branching & 0.13 & 2.87 & 2.87 & 0.955 & -0.08 & 0.74 & 0.74 & 0.940 \\
 &  &  & fitted fixed form & 0.02 & 3.29 & 3.29 & 0.960 & -0.04 & 0.84 & 0.84 & 0.950 \\
 &  &  & EM, fitted forms & 0.03 & 2.75 & 2.75 & 0.953 & 0.02 & 0.68 & 0.68 & 0.930 \\
 &  &  & EM, uniform forms & -0.46 & 4.16 & 4.19 & 0.963 & -0.01 & 1.00 & 1.00 & 0.938 \\
education & 4 & 5 & branching & -0.15 & 2.85 & 2.86 & 0.970 & 0.01 & 0.77 & 0.77 & 0.953 \\
 &  &  & fitted fixed form & -0.25 & 3.42 & 3.42 & 0.973 & 0.01 & 0.86 & 0.86 & 0.958 \\
 &  &  & EM, fitted forms & -0.63 & 2.81 & 2.88 & 0.968 & -0.45 & 0.74 & 0.87 & 0.920 \\
 &  &  & EM, uniform forms & -0.84 & 5.00 & 5.06 & 0.958 & -0.44 & 1.25 & 1.32 & 0.925 \\
education & 5 & 6 & branching & -0.15 & 2.83 & 2.83 & 0.970 & 0.01 & 0.77 & 0.77 & 0.948 \\
 &  &  & fitted fixed form & -0.04 & 3.02 & 3.01 & 0.978 & -0.00 & 0.86 & 0.86 & 0.938 \\
 &  &  & EM, fitted forms & -0.36 & 2.51 & 2.54 & 0.963 & -0.17 & 0.72 & 0.74 & 0.953 \\
 &  &  & EM, uniform forms & -0.65 & 4.14 & 4.18 & 0.978 & -0.20 & 1.03 & 1.05 & 0.950 \\
cost barrier & 4 & 5 & branching & 0.13 & 3.18 & 3.18 & 0.935 & -0.04 & 0.74 & 0.74 & 0.953 \\
 &  &  & fitted fixed form & -0.07 & 3.85 & 3.85 & 0.960 & -0.02 & 0.94 & 0.94 & 0.943 \\
 &  &  & EM, fitted forms & 0.79 & 2.99 & 3.09 & 0.920 & 0.68 & 0.70 & 0.98 & 0.870 \\
 &  &  & EM, uniform forms & 0.45 & 4.85 & 4.87 & 0.930 & 0.69 & 1.13 & 1.32 & 0.922 \\
cost barrier & 5 & 8 & branching & 0.01 & 3.08 & 3.08 & 0.938 & -0.03 & 0.74 & 0.74 & 0.953 \\
 &  &  & fitted fixed form & -0.11 & 3.79 & 3.78 & 0.943 & 0.01 & 0.93 & 0.93 & 0.953 \\
 &  &  & EM, fitted forms & 0.48 & 3.01 & 3.04 & 0.932 & 0.33 & 0.68 & 0.75 & 0.935 \\
 &  &  & EM, uniform forms & 0.41 & 3.71 & 3.73 & 0.963 & 0.45 & 0.89 & 0.99 & 0.945 \\
poor health & 4 & 4 & branching & -0.21 & 3.21 & 3.21 & 0.948 & -0.00 & 0.77 & 0.77 & 0.958 \\
 &  &  & fitted fixed form & -0.19 & 3.23 & 3.23 & 0.960 & -0.00 & 0.79 & 0.79 & 0.955 \\
 &  &  & EM, fitted forms & -0.72 & 3.07 & 3.15 & 0.943 & -0.60 & 0.74 & 0.95 & 0.885 \\
 &  &  & EM, uniform forms & -0.75 & 4.84 & 4.90 & 0.958 & -0.69 & 1.19 & 1.38 & 0.940 \\
poor health & 5 & 7 & branching & -0.28 & 3.26 & 3.27 & 0.948 & -0.03 & 0.77 & 0.76 & 0.958 \\
 &  &  & fitted fixed form & -0.37 & 4.00 & 4.02 & 0.953 & -0.07 & 0.99 & 0.99 & 0.948 \\
 &  &  & EM, fitted forms & -0.83 & 2.97 & 3.08 & 0.950 & -0.69 & 0.74 & 1.01 & 0.882 \\
 &  &  & EM, uniform forms & -0.73 & 4.20 & 4.26 & 0.948 & -0.62 & 1.05 & 1.21 & 0.890 \\
\bottomrule
\end{tabular}
\end{table}

\nEmAppendixText

\paragraph{Accuracy against audit size.}
For all \nCurveTrees{} BRFSS trees at caps four and five (two targets, three seeds) and audit sizes from 500 to 32,000, the design-based form score and the working models with all two-way and all three-way interactions are applied to the same fitted-form records (\nCurveReps{} replications per tree and size; no intervals).
\begin{table}[h]
\centering
\caption{Working models against the design-based form score on the same fitted-form records, as a function of the audit size $n$ (12 BRFSS trees: caps 4 and 5, two targets, three seeds; 400 replications per tree and size). Entries are medians over trees [minimum, maximum] of the RMSE relative to the design-based form score; EM-2 and EM-3 are log-linear models with all two-way and all three-way interactions. Worse: number of trees for which the model has the larger RMSE. Bias/SD: median absolute bias of the model in units of its own standard deviation.}
\label{tab:curve}
\scriptsize
\setlength{\tabcolsep}{3pt}
\begin{tabular}{rlrrlrrl}
\toprule
$n$ & EM-2 / form score & worse & bias/SD & EM-3 / form score & worse & bias/SD & branching / form score \\
\midrule
500 & 0.81 [0.77, 0.94] & 0/12 & 0.11 & 1.00 [0.95, 1.06] & 7/12 & 0.05 & 0.82 [0.77, 0.99] \\
1,000 & 0.81 [0.73, 0.93] & 0/12 & 0.08 & 0.98 [0.93, 1.00] & 1/12 & 0.05 & 0.82 [0.73, 0.98] \\
2,000 & 0.80 [0.70, 0.92] & 0/12 & 0.21 & 0.96 [0.88, 0.98] & 0/12 & 0.09 & 0.83 [0.74, 0.98] \\
4,000 & 0.82 [0.72, 1.00] & 0/12 & 0.21 & 0.94 [0.90, 1.01] & 1/12 & 0.15 & 0.80 [0.73, 0.98] \\
8,000 & 0.83 [0.70, 1.09] & 3/12 & 0.25 & 0.95 [0.89, 1.04] & 1/12 & 0.23 & 0.84 [0.73, 0.98] \\
16,000 & 0.83 [0.74, 1.36] & 4/12 & 0.34 & 0.96 [0.91, 1.10] & 3/12 & 0.37 & 0.82 [0.74, 0.98] \\
32,000 & 0.88 [0.79, 1.81] & 4/12 & 0.45 & 1.04 [0.90, 1.13] & 9/12 & 0.55 & 0.81 [0.74, 0.97] \\
\bottomrule
\end{tabular}
\end{table}

The two-way model has the smaller RMSE for every tree at $n\le\nCurveLastAllBetter$ and the larger RMSE for \nCurveTwoWorsec{} and \nCurveTwoWorsed{} of the \nCurveTrees{} trees at $n=16{,}000$ and $32{,}000$. Its median RMSE ratio stays between \nCurveTwoMedLo{} and \nCurveTwoMedHi{} at every size, so for two thirds of the trees the model keeps its advantage up to 32,000 records; which trees these are cannot be read from the capped records. The three-way model stays within \nCurveThreeLo--\nCurveThreeHi{} of the design-based score at every size and is the less accurate for \nCurveThreeWorsed{} of the trees at 32,000. With all interactions the model is saturated and its maximum-likelihood estimate is no longer a parsimonious alternative. A model-assisted score, which would use a working model only to choose the tables of a pointwise-unbiased score, cannot have a smaller variance than the oracle tables for the same form distribution, so it cannot recover the working model's advantage at small $n$; that advantage is bought with the model's bias.

\subsection{When Asking Changes the Answer}\label{app:measurement}
Let $P$ describe latent complete records and let $K_T$ be the distribution of the recorded transcript given a latent record under form $T$. A frozen score then estimates $\E_{P,q,K}H_T(O_T)$, which can differ from $\Delta$; a branching policy has its own sequential kernel. We compute the resulting biases exactly on the screening trees of Appendix~\ref{app:screening}. The answer at query position $j=1,\ldots,b$ flips independently with probability $\epsilon(j-1)/(b-1)$, so the first answer is exact and the last flips with probability $\epsilon$; labels are unchanged. Branching routes on the recorded answers; fixed forms ask their coordinates in increasing index order and use the fitted design's scores. Summing over all $2^b$ flip patterns gives the exact mean and second moment of each design at the target law. This is an illustrative position model, not an estimated response-error process.
\begin{table}[h]
\centering
\caption{Position effects on the screening trees: median absolute bias in units of the clean branching standard error at $n=1000$, and the share of runs in which branching keeps the smaller RMSE at $n=1000$ (the last answer flips with probability $\epsilon$).}
\label{tab:position}
\small
\begin{tabular}{llrrrr}
\toprule
Data & $\epsilon$ & branching & fitted form & oracle form & branching RMSE smaller \\
\midrule
BRFSS & 0.01 & 0.15 & 0.21 & 0.22 & 24/24 \\
BRFSS & 0.03 & 0.45 & 0.64 & 0.65 & 24/24 \\
BRFSS & 0.05 & 0.75 & 1.08 & 1.09 & 24/24 \\
BRFSS & 0.1 & 1.54 & 2.17 & 2.20 & 24/24 \\
Thyroid & 0.01 & 0.02 & 0.24 & 0.23 & 24/24 \\
Thyroid & 0.03 & 0.06 & 0.71 & 0.70 & 24/24 \\
Thyroid & 0.05 & 0.11 & 1.18 & 1.17 & 24/24 \\
Thyroid & 0.1 & 0.21 & 2.37 & 2.34 & 23/24 \\
\bottomrule
\end{tabular}
\end{table}


\subsection{Cost Regimes Under a Relaxable Cap}\label{app:costchoice}
Suppose each feature costs one unit and each label $c$ units, and the cap may be relaxed. A complete record costs $c+m$ and attains the full-information variance $\sigma_P^2$; a capped record costs at most $c+b$. For a screening tree, branching attains $\sigma_P^2$ at cost $c+b$ and therefore dominates complete records. Complete records are cheaper than the best fixed design exactly when $(c+m)/(c+b)<W_{\rm fixed}/\sigma_P^2$, that is, when $c>(m-\zeta b)/(\zeta-1)$ with $\zeta=W_{\rm fixed}/\sigma_P^2$. For the BRFSS trees ($m=12$) the median $\zeta$ is 1.49, 1.35, 1.58 at $b=4, 5, 6$, so complete records undercut the best fixed design once a label costs more than about 12, 15, 4 feature questions, respectively, while branching is never undercut. Under an immutable cap the complete-record alternative is unavailable.
